\section{Benchmark Results}
\label{sec:results}

Tables~\ref{tab:marginals}--\ref{tab:pairwise} and the masking diagnostics in Table~\ref{tab:masking} address three related questions: how input representation affects masked-token prediction, whether one encoder remains stronger under matched conditions, and whether that ranking holds when evaluation places greater emphasis on technical terminology. The primary results come from 175 configurations formed by crossing seven encoders, five representations, and five informativeness subsets. Each configuration uses the shared evaluation design described in Section~5, so differences can be examined without changing the underlying record sample.

\subsection{Representation and Subset Effects}
\label{sec:repr}
\input{tables/table3_marginals}

Table~\ref{tab:marginals} shows that representation has a marked association with Top-1 accuracy. The mixed representation has the highest marginal mean (47.13\%), followed by key--value (38.39\%), template (30.93\%), narrative (30.33\%), and JSON (22.04\%). Because these formats are derived from the same occurrence records, the spread is not attributable to differences in source incidents. It reflects how the format presents the information available to the model when a token is masked. In this benchmark, combining structured fields with prose provides the most favourable prediction context, whereas the JSON representation produces the lowest marginal accuracy.

Differences among informativeness subsets are smaller. The high subset has a mean Top-1 accuracy of 35.53\%, compared with 31.55\% for the low subset; random, medium, and balanced subsets fall between these values (34.59\%, 34.36\%, and 32.80\%, respectively). The 25.09-percentage-point spread across representations is 6.30 times the 3.98-point spread across subsets. Thus, the amount of domain information represented by a subset does not account for the larger variation associated with serialization.

This distinction is important when interpreting encoder scores. A score obtained on mixed text is not directly interchangeable with one obtained on JSON, even when the records are identical. Model comparisons are therefore most informative when the representation and subset are held constant, with marginal results used to describe the broader pattern rather than to replace matched comparisons.

\subsection{Encoder Ranking and Robustness}
\label{sec:ranking}
\input{tables/table4_leaderboard}

Across the 25 representation--subset conditions, ModernBERT has the highest mean maritime Top-1 accuracy at 57.99\%. RoBERTa is second at 46.12\%; the remaining encoders have lower aggregate scores (Table~\ref{tab:leaderboard}). Top-1 accuracy measures the proportion of masked target positions at which the correct token is the model's highest-ranked prediction. ModernBERT's score is a macro-average over the 25 conditions, not the result of one favourable cell. Its score in the JSON $\times$ Balanced condition is 58.14\%, illustrating the distinction between a single condition and the overall average.

The matched comparison with RoBERTa helps show how stable that advantage is. ModernBERT leads in 20 of the 25 conditions. Its only losing block is the narrative representation, where the representation-level means are 45.16\% for ModernBERT and 45.44\% for RoBERTa. The mean difference across all matched conditions is 11.87 percentage points. The reported bootstrap intervals, $52.81$--$63.42$ for ModernBERT and $40.92$--$51.01$ for RoBERTa, remain separated. The evidence therefore favours ModernBERT overall, while the narrative results show that the ranking is not identical for every format.

The composite score provides a similar ordering, with MECS values of 87.40 for ModernBERT and 87.04 for RoBERTa. These values should be interpreted in light of the scoring procedure defined in Section~5: the Top-1, rare-term, and loss components reach their upper normalized value for both leading models, leaving tokenizer fragmentation to account for the difference between their composite scores. MECS is therefore an operational screening score, not an independent confirmation of capability.

The multi-criteria analysis across six objectives identifies ModernBERT, RoBERTa, SciBERT, BERT-base, and Legal-BERT as non-dominated. A non-dominated model remains on the Pareto frontier because no other candidate is at least as good on all selected objectives and strictly better on one. The frontier preserves alternatives for settings where predictive performance must be balanced against vocabulary behaviour or inference cost. BERT-base has 110M parameters compared with 149M for ModernBERT and a measured latency of 19.51\,ms. SciBERT and Legal-BERT also remain non-dominated under the evaluated objectives, whereas BioBERT and PubMedBERT are dominated within this cohort.

The Top-1 difference between the highest- and lowest-ranked encoders is 37.70 percentage points, from ModernBERT at 57.99\% to PubMedBERT at 20.29\%. This ranking is specific to the present protocol and checkpoints. Since the encoders differ in pretraining data, scale, architecture, tokenizer, and training history, their relative performance cannot be interpreted as a controlled test of whether domain-specialized pretraining is intrinsically less effective.

\subsection{Tokenizer Coverage and Tokenization Dynamics}
\label{sec:tokenizer}

The tokenizer results help explain why vocabulary fit and contextual prediction do not produce the same ranking. BERT-base has the highest single-token coverage at 73.43\% and the lowest fragmentation rate among the compared encoders at 26.57\%, yet ranks fourth in mean Top-1 accuracy. ModernBERT has 37.01\% single-token coverage and 62.99\% fragmentation, but leads the primary benchmark (Table~\ref{tab:leaderboard}; Figure~\ref{fig:repr_tok}b). RoBERTa shows a similar coverage--fragmentation profile to ModernBERT, with 35.22\% single-token coverage and 64.78\% fragmentation.

Single-token coverage records whether a technical term is represented as one token; it does not measure whether the encoder can predict that token from context. A tokenizer may preserve a term as one unit while its model assigns that token a low rank in a sentence. Conversely, a model that divides a term into several pieces can still predict those pieces effectively. The results therefore do not support selecting an encoder on coverage alone. Under the primary random-masking protocol, ModernBERT achieves the strongest aggregate recovery despite its relatively high fragmentation. This conclusion is limited to the evaluated cohort and masking procedure; other tokenization schemes or downstream tasks may favour different trade-offs.

\begin{figure}[!htbp]
\centering
\vspace{2pt}
\includegraphics[width=\linewidth]{figures/fig2_representation_tokenizer_compact.pdf}
\vspace{4pt}
\caption{(a) Mean Top-1 accuracy (\%) across encoders and representations. (b) Single-token vocabulary coverage versus mean Top-1 accuracy. The panels show the relationship between input serialization, tokenizer coverage, and contextual prediction performance.}
\label{fig:repr_tok}
\end{figure}

\subsection{Crossed Factorial Statistical Validation}
\label{sec:stats}

The crossed repeated-measures analysis tests whether the differences observed in the marginal means are consistent across the matched design. Model identity accounts for 55.66\% of total variance and representation accounts for 25.47\%. Their interaction contributes a further 15.34\%, indicating that the effect of serialization depends on the encoder being evaluated. Subset accounts for 0.71\%; the Model $\times$ Subset and Representation $\times$ Subset interactions account for 0.59\% and 1.12\%, respectively. The residual pooled three-way interaction accounts for 1.11\% of variance ($\mathrm{df}=96$). These estimates put the relatively small subset contribution in context: subset differences are measurable, but model identity and representation explain substantially more of the variation in this experiment.

The ANOVA results are significant for Model, $F(6,96)=804.16$, $p<10^{-15}$, $\eta_p^2=0.9805$; Representation, $F(4,96)=552.03$, $p<10^{-15}$, $\eta_p^2=0.9583$; Model $\times$ Representation, $F(24,96)=55.40$, $p<10^{-15}$, $\eta_p^2=0.9327$; Representation $\times$ Subset, $F(16,96)=6.06$, $p=5.25\times10^{-9}$, $\eta_p^2=0.5026$; Subset, $F(4,96)=15.47$, $p=8.46\times10^{-10}$, $\eta_p^2=0.3919$; and Model $\times$ Subset, $F(24,96)=2.12$, $p=0.00533$, $\eta_p^2=0.3469$. The partial eta-squared values describe effect size within the fitted analysis; they should be read alongside, rather than substituted for, the variance contributions above.

The Friedman test provides a rank-based check on the overall model differences across the 25 matched conditions. It rejects the hypothesis that the seven encoders perform equivalently in this benchmark ($\chi^2=121.17$, $\mathrm{df}=6$, $p=9.27\times10^{-24}$, Kendall's $W=0.8078$). Kendall's $W$ indicates substantial agreement in the model rankings across conditions. This omnibus result establishes that the rankings differ; the pairwise comparisons show where the strongest separation occurs.

\input{tables/table5_pairwise}

After Holm correction, 19 of the 21 pairwise comparisons remain statistically significant. The clearest comparison is ModernBERT against RoBERTa: the mean difference is 11.87 percentage points, with an unadjusted 95\% confidence interval of $6.44$--$17.30$, an adjusted $p$-value of $2.29\times10^{-4}$, paired $d_z=0.9026$, and 20 wins in 25 matched conditions. The paired rank-biserial correlation is $r_{\mathrm{pb}}=0.8585$. Cliff's $\delta=0.5104$ describes a separate, unpaired comparison over 625 cross-pairs; it is not another estimate of the 20/25 matched-condition win rate. Against the other encoders, ModernBERT wins at least 24 of 25 conditions, with the reported $r_{\mathrm{pb}}\ge0.99$ and $\delta\ge0.83$.

The two comparisons that do not survive Holm correction are SciBERT--Legal-BERT ($\Delta=+4.42$ percentage points, $p_{\mathrm{adj}}=0.170$) and BERT-base--Legal-BERT ($\Delta=+0.11$ percentage points, $p_{\mathrm{adj}}=0.692$). These results do not establish equivalence; they indicate that the available matched-condition evidence does not distinguish the pairs after adjustment. The statistical tests support ModernBERT's overall lead within the evaluated design, but they do not establish that the same ranking will hold for every maritime corpus or downstream task.

\subsection{Masking Sensitivity and Pseudo-Log-Likelihood}
\label{sec:masking}
\input{tables/table6_masking}

The primary benchmark masks 15\% of tokenizer-defined input units at random. Table~\ref{tab:masking} reports a separate diagnostic in which domain-aware masking gives technical terms priority within the same masking budget. This diagnostic covers 28 model--condition configurations across four predefined representation--subset conditions per encoder, so its results should not be combined directly with the 175-cell leaderboard.

Domain-aware masking lowers Top-1 accuracy for every encoder in these diagnostic conditions. For example, ModernBERT falls from 62.44\% under random masking to 37.50\% under domain-aware masking, a decrease of 24.94 percentage points; RoBERTa falls from 50.08\% to 24.47\%, a decrease of 25.61 points. The rare-term results differ even more sharply from the primary ranking. Under random masking, rare-target Top-1 is 82.05\% for ModernBERT, 87.39\% for RoBERTa, and 48.18\% for SciBERT. Under DA-Rare, the corresponding values are 8.01\%, 53.23\%, and 38.71\%. ModernBERT's sharp drop shows that strong aggregate performance under random masking does not guarantee reliable recovery when the target distribution is concentrated on rare technical terms.

Sampled pseudo-log-likelihood (PLL) provides a complementary assessment of conditional token likelihood on sampled text. The resulting pseudo-perplexity is 13.72 for ModernBERT, 16.02 for RoBERTa, 25.50 for BERT-base, and 91.56 for PubMedBERT; lower values indicate better likelihood under this scoring procedure. These figures are auxiliary diagnostics, not substitutes for the primary MLM loss or perplexity. Token-level means also differ by token type: domain-token PLL is $-1.64$ for RoBERTa and $-1.95$ for ModernBERT, whereas general-token PLL is $-2.59$ for ModernBERT and $-2.83$ for RoBERTa. The less negative value is higher in each comparison. The diagnostics therefore give a more qualified picture than a single aggregate score: ModernBERT leads on sampled pseudo-perplexity, while RoBERTa has the higher reported mean domain-token PLL, and rare-term recovery remains sensitive to the masking protocol.

\subsection{Initialization Decision}
\label{sec:decision}

ModernBERT-base was selected to initialize domain-adaptive pretraining. It has the highest primary mean Top-1 accuracy (57.99\%), the leading reported MECS (87.40), 100\% bootstrap rank-1 stability, and no statistically significant pairwise defeats after Holm correction. Its lead over RoBERTa is supported across the matched design, although the 20 wins in 25 conditions and the narrative-format exception show that the advantage is not uniform across every condition. The decision is based on the combined ranking, matched comparisons, and Pareto analysis rather than MECS alone.

BERT-base remains a practical alternative where model size and latency are more important than maximizing the primary score: it has 110M parameters and a measured latency of 19.51\,ms. The rare-term diagnostic also gives a clear direction for follow-up evaluation. Domain adaptation should track performance on rare maritime terms explicitly rather than assume that an encoder's aggregate masked-token score guarantees performance on that subset.

The selection is an initialization decision for the present intrinsic benchmark, not evidence of downstream superiority. Masked-token recovery and PLL measure predictive behaviour on text; they do not establish gains in named-entity recognition, incident classification, retrieval, or operational decision support. As discussed in Section~9, only ModernBERT underwent DAPT in this study, and downstream evaluation remains necessary to establish whether this selection improves performance on applied maritime tasks.
